% Completion of the corrected elementary-representation induction for Kamp III.2, Lemma 8.

\subsubsection{Elementary boundary formulas and the universal-family induction}
\label{sec:p09-kamp-iii2-lemma8-completion}

The boundary decomposition becomes an elementary interval construction once
its bracketing conditions have actual elementary presentations. We first
construct those presentations, then localize each remaining universal family
at its interval and combine the resulting interval conditions. The final
induction proves the original conclusion of Kamp's Lemma~8 over
\(\mathsf{Com}\), with the corrected bracketing conditions
\cite[ch.~III, pp.~85--92]{Kamp1968}.

\begin{hbdefinition}{def:p09-kamp-order-dual}{Order reversal of a formula or word}
For a formula \(\theta\), let \(\theta^{\mathrm d}\) be obtained
by replacing every order atom \(r<s\) by \(s<r\), retaining
unary atoms and every Boolean operation and quantifier. This is an
operation on formulas, not an additional connective. For a structure
\(\mathcal M\), write \(\mathcal M^{\mathrm{op}}\) for the
same domain and unary interpretations with the order reversed.
For a labelled word \(W=(\alpha_0,\beta_1,\ldots,\beta_n,\alpha_n)\),
put
\[
 W^{\mathrm d}\defeq
 (\alpha_n^{\mathrm d},\beta_n^{\mathrm d},\alpha_{n-1}^{\mathrm d},
                    \ldots,\beta_1^{\mathrm d},\alpha_0^{\mathrm d}).
\]
Each label is reversed as a formula on its free parameter. Bound variables
are made fresh in every subsequent construction.
\end{hbdefinition}

\begin{hblemma}{thm:p09-kamp-elementary-duality}{proved=annex:p9.kamp.elementary.duality}{Order reversal preserves the elementary constructions}
Every \(\theta\in\operatorname{E}(w)\) has an explicitly
constructed \(\operatorname{E}(w)\) presentation of
\(\theta^{\mathrm d}\). Given an elementary word presentation
of \(\rho(x,y)\), reversing its word and its labels gives an
elementary presentation, using the constructed elementary representatives
for its dual labels, satisfying
\[
 \mathcal M\models\mathcal B_p(W^{\mathrm d};b,a)
 \quad\Longleftrightarrow\quad
 \mathcal M^{\mathrm{op}}\models\mathcal B_p(W;a,b)
 \qquad(b<a\text{ in }\mathcal M).
\]
Consequently a supplied property \(I(x,y)\) representation can be
reversed into an \(I(y,x)\) representation of its order dual.
Both \(\mathsf{Lin}\) and \(\mathsf{Com}\) are preserved
by order reversal. Basic words retain their occurrence lengths.
\end{hblemma}

\begin{hbdefinition}{def:p09-kamp-neighbor-labels}{One-sided neighborhoods as elementary labels}
For a unary letter or elementary label \(B\), define
\[
\begin{aligned}
 N_B^+(w)&\defeq\exists u\,(w<u\land\forall t\,(w<t<u\to B(t))),\\
 N_B^-(w)&\defeq\exists u\,(u<w\land\forall t\,(u<t<w\to B(t))).
\end{aligned}
\]
Use \(T^K\) for \(B\) to obtain the strict existence of a
later or earlier point. Each formula has an actual
\(\operatorname{E}(w)\) presentation: take the guarded future
or past closure of the one-interval elementary word \((T^K,B,T^K)\).

For a basic word \(W=(A_0,B_1,\ldots,B_n,A_n)\), keep the
limiting labels \(\lambda_W^+,\lambda_W^-\) of
\cref{def:p09-kamp-fixed-endpoint-limits}, and define
\[
\begin{aligned}
 e_W^+&\defeq\neg N_{T^K}^+\lor
                  (A_0\land N_{B_1}^+\land\lambda_W^+),\\
 e_W^-&\defeq\neg N_{T^K}^-\lor
                  (A_n\land N_{B_n}^-\land\lambda_W^-).
\end{aligned}
\]
All labels have the same free parameter \(w\). Their displayed
formulas stand for the corresponding explicit elementary constructions.
\end{hbdefinition}

\begin{hblemma}{thm:p09-kamp-unary-accumulation}{proved=annex:p9.kamp.unary.accumulation}{Global elementary labels for fixed-endpoint accumulation}
The labels \(e_W^+,e_W^-\) belong to \(\operatorname{E}\)
and, over \(\mathsf{Lin}\), express respectively
\(C_W^+\) and \(C_W^-\) from
\cref{def:p09-kamp-fixed-endpoint-limits}. The negative-neighborhood
disjunct handles a last or first point, where the defining universal
accumulation condition is vacuously true. This includes every
\(n\geq1\) and the separate one-letter limiting case.
\end{hblemma}

\begin{hbdefinition}{def:p09-kamp-remote-limits}{Realizations approaching the remote endpoint}
For a basic word \(W\), put
\[
\begin{aligned}
 B_W^+(x,y)&\defeq\forall v\bigl(v<y\to
  \exists t\,(x<t<y\land v<t\land\mathcal B_p(W;x,t))\bigr),\\
 B_W^-(x,y)&\defeq\forall v\bigl(x<v\to
  \exists t\,(x<t<y\land t<v\land\mathcal B_p(W;t,y))\bigr).
\end{aligned}
\]
Under \(x<y\), the instances \(v=x\) and \(v=y\)
respectively force an actual realization. These definitions are therefore
not vacuous when the outer endpoints are adjacent.

For \(1\leq k\leq n\), let \(p_{W,k}^-(w)\) be the
limiting-past formula for the suffix beginning at \(A_k\).
It is \(\partial^-A_n(w)\) when \(k=n\); otherwise it is
the \(\operatorname{Lp}_{n-k}(w)\) formula with word
\((A_k,B_{k+1},\ldots,B_n,A_n)\).
Let \(p_{W,k}^+(w)\) be the limiting-future formula for the
prefix ending at \(A_{k-1}\): it is \(\partial^+A_0(w)\)
when \(k=1\), and otherwise the corresponding
\(\operatorname{Lf}_{k-1}(w)\) formula. Define elementary words
\[
\begin{aligned}
 V_{W,k}^+&\defeq(A_0,B_1,\ldots,B_{k-1},A_{k-1},B_k,p_{W,k}^-),\\
 V_{W,k}^-&\defeq(p_{W,k}^+,B_k,A_k,B_{k+1},\ldots,B_n,A_n).
\end{aligned}
\]
The first word has \(k\) intervals; the second has \(n-k+1\).
The limiting labels are evaluated at the endpoint where they occur,
in the whole structure.
\end{hbdefinition}

\begin{hblemma}{thm:p09-kamp-remote-limits}{proved=annex:p9.kamp.remote.limits}{The complete remote-boundary representations}
Over \(\mathsf{Lin}\), under the guard \(x<y\),
\[
\begin{aligned}
 B_W^+(x,y)&\longleftrightarrow
             \bigvee_{k=1}^n\mathcal B_p(V_{W,k}^+;x,y),\\
 B_W^-(x,y)&\longleftrightarrow
             \bigvee_{k=1}^n\mathcal B_p(V_{W,k}^-;x,y).
\end{aligned}
\]
In particular both conditions have property \(I(x,y)\).
For a fixed initial endpoint, the first condition supplies the
remote cofinal boundary used as \(\beta_1\) in the source.
Its explicit \(x<t\) guard agrees with that source condition
under \(x<y\), as proved in the annex. There is no assumption
that different realizing sequences share their intermediate points
\cite[ch.~III, pp.~85, 89--90]{Kamp1968}.
\end{hblemma}

\begin{hbdefinition}{def:p09-kamp-interior-approach}{Approaching an interior point from beyond it}
For \(x<u\), define
\[
 A_W^+(x,u)\defeq\forall v\bigl(u<v\to
       \exists t\,(u<t<v\land\mathcal B_p(W;x,t))\bigr).
\]
When \(u\) has a later point, use every cut \(c\in\mathcal C_n\)
of \(W\) to form
\[
                    L_c(x,u)\land e_{W_c^R}^+(u).
\]
The second factor has an \(\operatorname{E}(u)\) presentation.
The finite disjunction of these factors has an \(I(x,u)\)
presentation by endpoint insertion and conjunction.
\end{hbdefinition}

\begin{hblemma}{thm:p09-kamp-interior-approach}{proved=annex:p9.kamp.interior.approach}{The unattained lower-boundary approach formula}
For every \(x<u\) in a linear structure,
\[
 A_W^+(x,u)\longleftrightarrow
 \left(\neg N_{T^K}^+(u)\ \lor\
   \bigvee_{c\in\mathcal C_n}
          (L_c(x,u)\land e_{W_c^R}^+(u))\right).
\]
The displayed formula has property \(I(x,u)\), including when
\(u\) is a last point. In an interior-boundary application
\(x<u<y\), the no-future disjunct is false.
Order reversal gives the corresponding fixed-final-endpoint formula
approaching an interior point from its left, with its explicit
elementary presentation. The finitely many cut alternatives are all
retained; one does not require the same cut at every initial witness.
\end{hblemma}

\begin{hblemma}{thm:p09-kamp-segmented-assembly}{proved=annex:p9.kamp.segmented.assembly}{Assembling interval conditions across finitely many ordered points}
Fix pairwise distinct variable names \(x_0,\ldots,x_m\), with \(m\geq1\),
and the strict order \(x_0<\cdots<x_m\).
Suppose finitely many conditions are given, each either an
\(\operatorname{E}(x_i)\) formula at a point, or a formula
with a supplied property \(I(x_i,x_j)\) representation for
some \(i<j\). Then the conjunction of those conditions,
existentially quantifying \(x_1,\ldots,x_{m-1}\) with the
strict order guard, has property \(I(x_0,x_m)\).

The finite construction expands every supplied interval disjunction,
cuts each elementary word at every intervening named point, distributes
over all cuts, fuses the resulting conditions on each adjacent interval,
and concatenates at the elementary point labels. Empty conjunctions use
the declared truth labels. Every imposed interval edge and every point
condition is retained. Input \(I\) equivalences retain their
\(\mathsf{Com}\) scope throughout this construction.
\end{hblemma}

\begin{hbdefinition}{def:p09-kamp-localization-packages}{Localizing a smaller family at one of the boundary intervals}
Let \(F\) have maximal basic member \(\phi\), write
\(D=D_{F\setminus\{\phi\}}\), and fix ordered endpoints
\(a\leq r<s\leq b\), where equality means that the endpoint
name is literally reused. Let \(\eta\in\operatorname{E}(t)\)
be a chosen additional current-point label, possibly \(F^K_t\).
The current-point variable \(t\) is distinct from all endpoint names.
Localize
\[
             \forall t\,(r<t<s\to D(t)\lor\eta(t))
\]
as follows. For an original \(\rho(a,t)\), retain it when
\(r=a\); otherwise apply all cut choices at \(r\), recording
fixed prefixes on \((a,r)\) and residuals on \((r,t)\).
For an original \(\rho(t,b)\), retain it when \(s=b\);
otherwise record all fixed suffix choices on \((s,b)\) and
residuals on \((t,s)\). Take every tuple of choices for every
crossing input. Retain the original \(\operatorname{At}(t)\)
members.

For each tuple, let \(G\) be the resulting local residual family
together with one fresh atom \(Q(t)\). It belongs to
\(\operatorname{At}(t)\cup\operatorname{Bp}(r,t)\cup
\operatorname{Bp}(t,s)\). The original member \(\phi\) is
absent. A \emph{localization package} consists of all selected
fixed prefix/suffix conjuncts and the local universal formula for
\(G\), properly substituting \(\eta\) for \(Q\).
The fresh letter occurs nowhere in the input or in \(\eta\).
\end{hbdefinition}

\begin{hblemma}{thm:p09-kamp-localization}{proved=annex:p9.kamp.localization}{Localization with an actual smaller-family induction hypothesis}
Assume the Lemma~8 conclusion for every finite family \(G\prec F\)
and every choice of its three endpoint/current-point variables.
Then the preceding localization construction gives a finite disjunction
of packages equivalent over \(\mathsf{Com}\) to the indicated
universal formula on \((r,s)\). Each package contains finitely
many basic fixed-side interval conditions and an actual
\(\operatorname{Ep}(r,s)\) disjunction obtained from the
induction hypothesis and proper substitution.

Every induction call has \(G\prec F\): deletion removes one
maximal member, each other original basic member contributes at most
one residual of its length and none longer, and the fresh atom has
length one. Equal outputs can only decrease the set count. The
elementary label \(\eta\) is inserted after that call; it is
not reclassified as an \(\operatorname{At}\) formula.
\end{hblemma}

\begin{hbdefinition}{def:p09-kamp-positive-boundaries}{Elementary support conditions for the eleven cases}
First take the left orientation \(H(t)=\mathcal B_p(W;a,t)\),
and let \(P(t)\) be its guarded past closure. Use the regional
matrices \(K_i\) and endpoint tests of
\cref{def:p09-kamp-corrected-boundary-branches}, but replace the
full descriptor \(\Gamma_i\) by the following positive supports.
For a two-boundary case with lower type \(\lambda\) and upper
type \(\mu\), the lower supports are
\[
 \begin{array}{c|l}
 0&l=a,\quad e_W^+(a),\\
 1&a<l,\quad\mathcal B_p(W;a,l),\\
 2&a<l,\quad A_W^+(a,l),
 \end{array}
\]
and the upper supports are
\[
 \begin{array}{c|l}
 0&h=b,\quad B_W^+(a,b),\\
 1&h<b,\quad\mathcal B_p(W;a,h),\\
 2&h<b,\quad B_W^+(a,h).
 \end{array}
\]
Retain \(a\leq l<h\leq b\), all outer-region \(D\) tests,
all middle-region \(D\lor P\) tests, and the \(D\) test at
each type-2 interior boundary. Quantify the same interior bounds
as in the eleven-case construction. Call the resulting formulas
\(Y_i\), for \(i=2,3,4,6,7,8,9,10,11\).
Set
\[
\begin{aligned}
 Y_1&\defeq\mathcal V_D(a,b),\\
 Y_5&\defeq\exists c\,(a<c<b\land\mathcal B_p(W;a,c)
                    \land\mathcal V_D(a,c)\land\mathcal V_D(c,b)).
\end{aligned}
\]
These supports retain actual fixed-start witnesses on both sides of
every middle point. They omit absence assertions only where the
unchanged \(D\) conditions already suffice for the target.
Their relation to the full descriptors is proved below, rather than
assumed as an identity of printed source formulas.
\end{hbdefinition}

\begin{hblemma}{thm:p09-kamp-positive-boundaries}{proved=annex:p9.kamp.positive.boundaries}{Correctness of the elementary-supported branch replacement}
For the left orientation, over \(\mathsf{Com}\) and under
\(a<b\),
\[
 \forall t\,(a<t<b\to D(t)\lor H(t))
                  \longleftrightarrow\bigvee_{i=1}^{11}Y_i.
\]
Each positive support has an actual \(\operatorname{E}\) or
\(I\) presentation as supplied above. This assertion does not
discard a bracketing hypothesis: the lower support supplies a
fixed-start witness strictly below every middle point, and the
upper support one strictly above it. Each \(Y_i\) individually
implies the target on every linear order. Complete orders supply
the full boundary descriptor needed to choose a branch in the
forward direction.
\end{hblemma}

\begin{hbmethod}{def:p09-kamp-lemma8-algorithm}{The terminating elementary-representation construction}
Given the exact finite family \(F\) of
\cref{def:p09-kamp-universal-family}, use its pair
\((M(F),c(F))\) for the recursion. For the point-only or empty
family, return the elementary word supplied by
\cref{thm:p09-kamp-family-normalization}. Otherwise choose a maximal
basic member. If it is right-oriented, first reverse the order
presentation and the basic words, perform the following left construction
with that distinguished reversed word, and reverse every resulting
elementary word back.

For the left construction form all eleven \(Y_i\). Replace the
guarded closure \(P\) by its actual \(\operatorname{E}\)
presentation from Lemma~1. For every regional universal condition
use the localization construction, with additional label \(P\)
in the middle regions and \(F^K\) in the outer regions. Recursively
obtain each smaller local family's representation, then perform the
specified proper substitution.

Expand every point condition \(D(l)\) or \(D(h)\) into its
original finite disjunction. Each alternative is either an elementary
point label or one original basic interval condition at its endpoints.
Expand every boundary support and every localization package into its
supplied interval disjunctions, distributing every finite choice.
Each resulting conjunction has only elementary point conditions and
interval edges between its at most four ordered named points.
Apply \cref{thm:p09-kamp-segmented-assembly} to each conjunction
and concatenate all returned disjunctions. No truth decision about the
structure is part of this syntax construction.
\end{hbmethod}

\begin{hblemma}{thm:p09-kamp-iii2-lemma8}{proved=annex:p9.kamp.iii2.lemma8}{Kamp's universal interior-disjunction lemma}
Let \(x,y,z\) be distinct variables and let
\[
 F\subseteq\operatorname{At}(z)\cup\operatorname{Bp}(x,z)
                                      \cup\operatorname{Bp}(z,y)
\]
be finite, including the empty family with the stated false disjunction.
Then
\[
                   \forall z\,(x<z<y\to\bigvee_{\psi\in F}\psi)
\]
has property \(I(x,y)\). The algorithm above returns a finite
disjunction of actual \(\operatorname{Ep}(x,y)\) formulas
equivalent to it under \(x<y\) over \(\mathsf{Com}\).
All original finite-family, variable and order hypotheses are retained.
The full proof includes both orientations and every one of the eleven
boundary cases, with strict decrease at every recursive call
\cite[ch.~III, lemma~8, pp.~83--92]{Kamp1968}.
\end{hblemma}

\begin{hbexample}{def:p09-kamp-remote-single-interval}{The terminal label is a strict limit}
For \(W=(A,B,C)\), the remote cofinal representative is
\[
                      (A,B,\partial^-C).
\]
Thus it requires \(A(a)\), \(B\) throughout \((a,b)\),
and \(C\)-points arbitrarily close below \(b\); it does not
require \(C(b)\). On \(\mathbb R\), take \(a=0,b=1\),
\(A=\{0\}\), \(B=\mathbb R\), and \(C=(0,1)\).
Realizations from 0 finish arbitrarily close below 1 although
\(C(1)\) is false. On the two-point order \(0<1\), the
remote condition is false even if every predicate is true, because
its instance at the initial endpoint requires a point strictly between
0 and 1. The elementary terminal limit is false there as well.
\end{hbexample}

\paragraph{The next source result.}
Lemma~9 begins on printed p.~92. It uses the completed universal-family
lemma to treat elementary point and interval conditions under a universal
interior quantifier. Its substitution and distribution construction is
the next separate result; Lemma~8 above no longer depends on it.
